\chapter{Current measurements with imperfect physical access}
\label{robustphysics:chapter}

The selective-channel current assay can tolerate quantified preparation,
control, and readout imperfections. Its useful extension is a matched native
control: preparing both arms from the same ensemble cancels a preparation
penalty that otherwise grows when the pulse is shortened. Common generator
drift also cancels to first order. Readout delay admits a conservative safe
window expressed through local endpoint hazards.

These results concern a declared finite classical Markov model and calibrated
error envelopes. They do not certify those envelopes from the response data,
and do not turn a sign crossing into an unconditional quantum witness. The
mathematical ingredients are standard semigroup perturbation, coupling, and
oscillation contraction; the contribution here is their explicit assembly into
a checkable physical current assay.

\section{The reference experiment}

Use row generators and row probability distributions. Let $Q$ be a finite
generator with stationary distribution $\pi$, and let a separately selectable
bidirectional physical channel between $i\ne j$ have rates $a,b\ge0$.
Its generator $L$ has $L_{ij}=a$, $L_{ji}=b$, $L_{ii}=-a$, $L_{jj}=-b$,
with all other entries zero. Set
\[
 s=1-r\in(0,1],\qquad Q_r=Q-sL,\qquad
 P(t)=e^{tQ},\quad P_r(t)=e^{tQ_r},\quad f=\mathbf1_j.
\]
Both $Q,Q_r$ must be valid generators. The native selected-channel current is
$J=\pi_i a-\pi_j b$. Multiple physical channels may connect the same states;
selectivity refers to the specified physical channel, not an unobserved choice
of algebraic decomposition.

For a stationary preparation, the native-minus-controlled contrast obeys
\begin{equation}\label{robustphysics:ideal}
 D_\pi(t)=\pi[P(t)-P_r(t)]f=sJH(t),\qquad
 t e^{-\Lambda_p t}\le H(t)\le t,
\end{equation}
where $\Lambda_p$ bounds the controlled exit rates of both endpoints. The
positive kernel is
$H(t)=\int_0^t(P_{r,jj}(u)-P_{r,ij}(u))\,du$. Its positivity follows from
the first-hitting argument in the finite-pulse theorem; its integrand need not
be positive. Thus the reference experiment does not require the surrounding
network to be reconstructed. Current and rates have units of inverse time;
$H$ and all pulse or delay durations have units of time. A response contrast
and all probability-error bounds are dimensionless.

\section{Generator error as a coupling certificate}

For a matrix $E$ with zero row sums define
\[
 \|E\|_{\mathrm{rowTV}}=\max_x\frac12\sum_y|E_{xy}|.
\]
The diagonal is included. For $E=A-B$, where $A,B$ are generators, this is
\[
 \max_x\max\left\{
  \sum_{y\ne x}(A_{xy}-B_{xy})_+,
  \sum_{y\ne x}(B_{xy}-A_{xy})_+\right\}.
\]

\begin{theorem}[Dimension-independent generator error]
\label{robustphysics:coupling}
Let $A(u),B(u)$ be bounded measurable finite generators on a common state
space, and start their processes from the same distribution. If
$\|A(u)-B(u)\|_{\mathrm{rowTV}}\le\kappa(u)$, then
\[
 \operatorname{TV}(\mathcal L(X_t),\mathcal L(Y_t))
 \le 1-\exp\left[-\int_0^t\kappa(u)\,du\right].
\]
\end{theorem}

\begin{proof}
While both processes are at $x$, couple their equal transitions to $y$ at
rate $\min(A_{xy},B_{xy})$. Let $U,V$ be the total remaining jump rates
of the two processes. Their remaining destinations are disjoint. Synchronize
remaining jumps to distinct destinations at rate $\min(U,V)$; run the excess
$|U-V|$ unilaterally. The total separation hazard is $\max(U,V)$, bounded
by $\kappa(u)$. The integrated hazard bounds the probability of ever separating
before $t$. Terminal disagreement and total variation are no larger. Complete
the coupling arbitrarily after separation.
\end{proof}

This rate can be half the looser sum of absolute off-diagonal discrepancies.
Rows $(-3,2,1)$ and $(-3,1,2)$ have row-TV discrepancy $1$, while that
off-diagonal sum is $2$. A one-way jump against a frozen process attains
$1-e^{-\kappa t}$, so no universally smaller bound follows from $\kappa$
alone. The certificate can cover a full piecewise intervention and readout
history, provided its integral covers that entire history.

\section{Matched preparation}

For a real column $g$, write
$\operatorname{osc}(g)=\max g-\min g$. A stochastic matrix contracts
oscillation, and a zero-mass signed row $v$ satisfies
$|vg|\le\operatorname{TV}(v)\operatorname{osc}(g)$, where
$\operatorname{TV}(v)=\tfrac12\sum_x|v_x|$.

\begin{theorem}[Preparation cancellation]
\label{robustphysics:preparation}
Suppose both arms use the same preparation $\rho$, with
$\operatorname{TV}(\rho,\pi)\le\eta$. Then
\begin{align}
 |D_\rho(t)-D_\pi(t)|
 &\le\min\left\{
 \eta s(a+b)t,
 2\eta\bigl(1-e^{-s\max(a,b)t}\bigr),2\eta\right\}.
 \label{robustphysics:preparationbound}
\end{align}
The same statement holds after any common stochastic terminal operation $R$,
with $f$ replaced by $Rf$.
\end{theorem}

\begin{proof}
Duhamel's identity gives
\[
 P(t)-P_r(t)=s\int_0^tP_r(u)L P(t-u)\,du.
\]
The only nonzero entries of $Lg$ are $a(g_j-g_i)$ at $i$ and
$-b(g_j-g_i)$ at $j$. Consequently
$\operatorname{osc}(Lg)=(a+b)|g_j-g_i|\le(a+b)\operatorname{osc}(g)$.
Oscillation contraction bounds that of $[P(t)-P_r(t)]f$ by $s(a+b)t$.
Apply the signed-row inequality to $\rho-\pi$. Alternatively, the coupling
rate for $Q,Q_r$ is $s\max(a,b)$, so each entry of this column has absolute
value at most $1-e^{-s\max(a,b)t}$. Its oscillation is at most twice that
quantity, and is always at most $2$. Finally $Rf$ remains in $[0,1]$.
\end{proof}

The coefficient $\eta s(a+b)$ is attained at first order when
$\rho-\pi=(\eta,-\eta)$ on the two endpoints. Dividing by $st$ to
estimate current gives a bounded preparation error $\eta(a+b)$ as $t$ tends
to zero. Using the older response bound $2\eta$ would instead produce
$2\eta/(st)$. The improvement requires actual matching of the preparation
distributions, not merely the same stated preparation method.

If the controlled and native arms use $\rho_c,\rho_n$, one may apply the
matched bound to $\rho_n$ and add
$\zeta=\operatorname{TV}(\rho_c,\rho_n)$ to the response error. That
unmatched term can again be amplified by $1/t$. Random assignment to a common
ensemble supports matching; it does not establish stationarity or independent
records under uncontrolled drift. A steady observed mean cannot by itself
calibrate the TV distance to a hidden stationary distribution.

\section{A safe window for readout delay}

After the pulse, suppose both arms evolve under the native generator for
delay $d$, so $R(d)=e^{dQ}$. Let $\lambda_j$ bound the native exit rate from
$j$, and let $\beta_j$ bound $\max_{x\ne j}Q_{xj}$. Define
\begin{equation}\label{robustphysics:omega}
 \omega(d)=\operatorname{osc}((R(d)-I)f)
 \le\bar\omega(d):=(1-e^{-\lambda_j d})+(1-e^{-\beta_j d}).
\end{equation}
Indeed, the $j$ entry is $P_{jj}(d)-1\le0$, with magnitude at most the
probability of a departure. Every other entry is $P_{xj}(d)\ge0$, bounded
by the probability of entrance to $j$. Before that entrance its hazard is
at most $\beta_j$. Their range gives \eqref{robustphysics:omega}.

\begin{theorem}[Delayed sign protection]
\label{robustphysics:delay}
The stationary delayed contrast is $D_\pi(t,d)=sJK(t,d)$, where
\[
 K(t,d)=\int_0^t(e_j-e_i)P_r(u)R(d)f\,du,
\]
and
\[
 \max\{-t,t(e^{-\Lambda_p t}-\bar\omega(d))\}
 \le K(t,d)\le t.
\]
In particular, $\bar\omega(d)<e^{-\Lambda_p t}$ guarantees the native
current sign despite the delay.
\end{theorem}

\begin{proof}
Multiplying the stationary response identity by $R(d)$ gives the expression
for $K$. Each integrand is a difference of probabilities, so $|K|\le t$.
Furthermore
\[
 |K-H|\le\int_0^t
 |(e_j-e_i)P_r(u)(R(d)-I)f|\,du\le t\omega(d).
\]
Combine this with \eqref{robustphysics:ideal}.
\end{proof}

This is a sufficient safe window. Failure of its inequality means that the
certificate abstains; it does not establish a reversal. Unrestricted readout
delay can reverse a classical response, as the earlier counterexample shows.

Now let a clock select a duration $T\in[t_-,t_+]$, $t_->0$, independently
of the hidden state and preparation, with the same law in both arms. Its
distribution need not be known. For a valid uniform delay bound $\bar\omega$,
the averaged kernel satisfies
\begin{equation}\label{robustphysics:jitter}
 L:=\max\left\{-t_+,\min_{u\in\{t_-,t_+\}}
 u(e^{-\Lambda_pu}-\bar\omega)\right\}
 \le\mathbb E K(T,d)\le t_+.
\end{equation}
To see why endpoints suffice, the derivative of
$u(e^{-\Lambda_pu}-\bar\omega)$ is
$e^{-\Lambda_pu}(1-\Lambda_pu)-\bar\omega$. Before $1/\Lambda_p$ its
positive part decreases strictly, and afterwards it is nonpositive. Thus its
only possible interior extremum is a maximum; for $\Lambda_p=0$ the function
is linear. With no delay, every positive duration has a positive kernel.
State-dependent stopping is a different experiment and is not covered.

\section{An exact current interval}

A common fixed binary detector has mean $o+gp_j$. The offset $o$ cancels
in contrasts. Assume $0<g_-\le g\le g_+\le1$; a known negative gain can
instead be handled by relabeling detector outputs. Arbitrary coarse partitions
are not affine singleton detectors. Setting-dependent detector changes require
their own contrast error bound $\epsilon_{\rm det}$.

Let $\eta$ bound common preparation error, $\zeta$ its mismatch between
arms, and $K_n,K_c$ the calibrated integrated row-TV generator errors of the
full native and controlled histories. With rate caps $a\le a_+$,
$b\le b_+$, put
\begin{align*}
 b_p&=\min\{\eta s(a_++b_+)t_+,
 2\eta(1-e^{-s\max(a_+,b_+)t_+}),2\eta\},\\
 b&=g_+\{b_p+\zeta+(1-e^{-K_n})+(1-e^{-K_c})\}
       +\epsilon_{\rm det}.
\end{align*}
If a sampling procedure supplies an interval $[D_-,D_+]$ for the actual
native-minus-controlled recorded mean, widen it by $b$ and intersect with
$[-g_+,g_+]$. If the lower kernel bound $L$ in
\eqref{robustphysics:jitter} is positive, the ideal mean is $cJ$ for some
\[
 c\in[g_-sL,g_+st_+].
\]
The four endpoint ratios enclose $J$. Intersect also with the physical cap
$[-b_+,a_+]$. An interval excluding zero certifies current orientation;
an empty interval indicates inconsistent data and calibration inputs. If
$L\le0$, this certificate does not establish a sign. Joint calibration
failure probability adds once to the sampling failure probability.

The implementation uses exact rational arithmetic and positive Taylor
remainder bounds for exponentials. It cannot verify an apparatus merely by
receiving its asserted calibration numbers. Systematic errors do not disappear
with additional records, and preparation, clocks, or generator errors cannot
be treated as common unless that premise is justified.

\section{Current measurement without stationary preparation}

When stationarity cannot be established, the prepared ensemble's instantaneous
current can still be a useful physical target. For an arbitrary common
preparation $\rho$, put $J_\rho=\rho_i a-\rho_j b$.

\begin{theorem}[Stationarity-free short-pulse assay]
\label{robustphysics:nonstationary}
Let $k_i',k_j'$ each bound both departure from and maximal entrance to the
corresponding endpoint under $Q_r$. Let $\lambda_j,\beta_j$ bound native
departure and entrance as above, and suppose $|J_\rho|\le J_+$.
Then
\begin{equation}\label{robustphysics:nonstationarybound}
 \left|\frac{D_\rho(t)}{st}-J_\rho\right|
 \le\frac t2\{a k_i'+b k_j'+J_+(\lambda_j+\beta_j)\}.
\end{equation}
No stationary distribution is required.
\end{theorem}

\begin{proof}
Duhamel gives the exact identity
\[
 D_\rho(t)=s\int_0^t J_{\rho P_r(u)}
                [P_{jj}(t-u)-P_{ij}(t-u)]\,du.
\]
State $x$ membership can change only by departing from or entering $x$.
The probability of such a change before $u$ is at most
$1-e^{-k_x'u}\le k_x'u$, from any starting distribution. Consequently
$|J_{\rho P_r(u)}-J_\rho|\le a k_i'u+b k_j'u$.
For $h(v)=P_{jj}(v)-P_{ij}(v)$, one has $|h(v)|\le1$ and
\[
 |h(v)-1|\le(1-e^{-\lambda_jv})+(1-e^{-\beta_jv})
             \le(\lambda_j+\beta_j)v.
\]
Subtract $J_\rho$ inside the integral and integrate the two bounds.
\end{proof}

Without a sharper input, $J_+=\max(a,b)$ is valid. This theorem targets
the instantaneous current of the prepared ensemble. It does not identify an
unspecified steady-state flux or extend the all-time sign theorem to arbitrary
preparations. It is a finite-difference assay with an explicit bias bound.

\section{Shared drift cancels to first order}

The generic error certificate treats the arms separately. A genuinely shared
additive generator error supports a stronger conclusion.

\begin{theorem}[Common generator drift]
\label{robustphysics:shareddrift}
Suppose the actual pulse generators are $Q+E$ and $Q_r+E$, all four
generators being valid, and $\|E\|_{\mathrm{rowTV}}\le\kappa$. With the
same arbitrary initial distribution and terminal effect $f\in[0,1]$, let
$D^E,D$ denote actual and reference matched contrasts. Then, with
$M=\max(a,b)$,
\begin{align*}
 |D^E-D|&\le s(a+b+2M)\int_0^t(1-e^{-\kappa u})\,du\\
         &\le \tfrac12 s(a+b+2M)\kappa t^2.
\end{align*}
\end{theorem}

\begin{proof}
Mark actual semigroups by a superscript $E$. Subtract the two Duhamel
response expressions to obtain
\begin{align*}
 D^E-D=s\int_0^t\rho\big[&
 (P_r^E(u)-P_r(u))L P^E(t-u)\\
 &+P_r(u)L(P^E(t-u)-P(t-u))\big]f\,du.
\end{align*}
The first row difference has TV at most $1-e^{-\kappa u}$ and the column
$LP^Ef$ has oscillation at most $a+b$. For the second term,
$\rho P_r(u)L$ is its instantaneous channel current, of absolute value at
most $M$, times $e_j-e_i$. The remaining column has oscillation at most
$2(1-e^{-\kappa(t-u)})$. Integration proves the first inequality and
$1-e^{-\kappa u}\le\kappa u$ proves the second.
\end{proof}

For $\kappa>0$ the exact integral is
$t-(1-e^{-\kappa t})/\kappa$; it is zero at $\kappa=0$. After division
by $st$, shared drift contributes at most
$(a+b+2M)\kappa t/2$ to current error, which vanishes as the pulse shrinks.
The same norm bound for two different errors is insufficient: their matrices
must be the same. A common terminal operation is allowed, but an actual versus
reference change of the postpulse delay semigroup needs its own bound. This
pulse-only improvement is provided as a separate helper so that it cannot be
silently substituted into the complete independent-arm certificate.

\section{Evidence and operational limits}

The exact module \texttt{robust\_physics.py} implements generator discrepancy,
matched preparation, delayed kernels, current intervals, short-pulse bias, and
shared-drift bounds. Fourteen tests cover the arithmetic and stated contracts.
An additional seeded numerical sweep checks all five principal bounds on
$180$ constructed positive networks with two through eight states; matrix
exponentials are numerical checks, not the proofs. The fixed seed is $71293$.
All generated records and models are synthetic. No empirical physical
validation is reported.

The main experimental improvement is a choice between two explicit targets:
native stationary current with a calibrated preparation envelope, or
instantaneous prepared-ensemble current with a short-pulse bias bound. Both
need matched preparation, selective control or a calibrated spillover bound,
and a justified singleton or protected detector. The delayed safe window and
shared-error cancellation are useful only with their corresponding physical
rate and matching certificates. The formulas identify the missing calibration
when inference fails; they do not replace that calibration.
